\chapter{Introduzione}

\section{Contesto biologico per informatici}
Le \glspl{protein} sono i mattoni fondamentali di una cellula e il suo DNA, il \gls{genome}, codifica le informazioni necessarie a costruirle. Ogni cellula di un organismo pluricellulare porta essenzialmente lo stesso DNA, eppure le cellule si sviluppano in tipi distinti. La differenza non può venire dal DNA stesso. Tra la lettura del genoma e la costruzione delle proteine si colloca quindi uno o più livelli di regolazione che determinano quali proteine una cellula produce, e sono proprio quelle proteine a conferirle le sue caratteristiche. Caratterizzare questi livelli è centrale per spiegare come funzionano le cellule e che cosa ne orienta il comportamento.

L'espressione genica procede in tre fasi. Un \gls{gene} viene dapprima copiato in una lunga molecola di \gls{rna}, detta \gls{unsplicedTranscript}. La seconda fase, lo \gls{splicing}, rimuove i blocchi non codificanti e unisce quelli codificanti nelle \glspl{spliceJunction} risultanti, producendo un \gls{transcript} spliced più corto. La versione finale, spliced, è l'\gls{mrna}, che il ribosoma traduce\defn{Traduzione}{Il processo con cui il ribosoma legge una molecola di mRNA e assembla la corrispondente catena di amminoacidi in una proteina.} in proteina nella terza e ultima fase.

Il processo di accorciamento non è una mappatura uno-a-uno. Per ampliare la capacità di decodifica di un gene, le cellule possono produrre diverse versioni spliced dello stesso \gls{unsplicedTranscript}; queste versioni sono dette \glspl{isoform}. Poiché isoforme diverse possono tradursi in proteine con ruoli del tutto diversi, o persino contrastanti, una cellula può alterare profondamente il proprio comportamento semplicemente cambiando la variante che esprime---l'isoforma MET$\Delta$14, per esempio, agisce da interruttore che guida la crescita invasiva nel cancro \cite{cerqua2022met}. L'obiettivo centrale dell'analisi dell'\gls{dtu} è quindi misurare i cambiamenti nelle proporzioni di espressione relativa di specifiche isoforme di trascritto dello stesso gene tra condizioni biologiche, stadi di sviluppo o tipi cellulari diversi, per portare alla luce eventi regolatori che la misura a livello genico rende invisibili.

\section{Breve storia del sequenziamento a singola cellula}
Per gran parte della sua storia, l'\gls{rnaSeq} è stato eseguito su campioni bulk. L'RNA viene estratto da una popolazione di cellule, riunito in un'unica reazione e sequenziato come una sola libreria\defn{Libreria di sequenziamento}{L'insieme di molecole di DNA caricate nello strumento per una corsa. Qui contiene l'RNA del campione, copiato in DNA ed etichettato in modo che lo strumento possa leggere ogni frammento.}, quindi il profilo risultante descrive il campione nel suo insieme e non le singole cellule che lo compongono. Il sequenziamento di questo tipo, detto \gls{bulkRnaSeq}, è statisticamente comodo: riunire da migliaia a milioni di cellule produce una copertura profonda e stime di abbondanza stabili per ogni gene. La media, però, è anche il suo limite principale. L'abbondanza misurata di un gene è la media su tutte le cellule presenti, quindi l'eterogeneità cellulare\defn{Eterogeneità cellulare}{La coesistenza di tipi e stati cellulari distinti all'interno di un singolo organismo.} svanisce prima che i dati raggiungano chi li analizza. Un tipo cellulare raro che costituisce il due per cento di un tessuto contribuisce per il due per cento delle \glspl{read} di quel tessuto per ciascuno dei suoi marcatori\defn{Geni marcatore}{Un gene o trascritto la cui espressione identifica specificamente un tipo cellulare, uno stato o una condizione.} ed è facilmente sepolto sotto la popolazione dominante; due sottopopolazioni che si spostano in direzioni opposte---una che acquista un trascritto mentre l'altra lo perde---si annullano a vicenda e appaiono invariate.

Lo \gls{scRnaSeq} ha sostituito quella visione aggregata con una profilazione per cellula, eseguendo la stessa misura in modo indipendente su migliaia di cellule. La strategia si fonda sulla compartimentazione. Ogni cellula viene isolata fisicamente, il suo RNA poliadenilato catturato e retrotrascritto separatamente, e ogni molecola risultante marcata con due etichette. Un \gls{cellBarcode} identifica la cellula di origine, così i read possono essere raggruppati di nuovo in profili per cellula dopo il sequenziamento; un \gls{umi} contrassegna la singola molecola, così i duplicati di \gls{pcr} vengono collassati invece di essere contati erroneamente come espressione. Una singola corsa di sequenziamento riunita viene così convertita in migliaia di misure di espressione indipendenti, ciascuna molto meno profonda di un campione bulk.

Il primo \gls{scRnaSeq}, prodotto dal 2009 in poi, proveniva da protocolli \gls{plateProtocols}, che consumano una cellula per reazione \cite{tang2009mrna}. La chimica di quell'epoca ancorava i read a un'estremità della molecola, e la precisione delle basi diminuisce lungo un read man mano che si allunga.\defn{Base}{Uno dei quattro nucleotidi che formano un filamento di DNA o RNA: adenina (A), citosina (C), guanina (G) e timina (T). Un read è una sequenza di basi consecutive, e la sicurezza con cui ogni base viene identificata diminuisce lungo il read man mano che si allunga.} L'informazione affidabile si concentrava quindi vicino a quell'estremità, un limite che un sequenziamento più profondo non poteva rimediare---i dati di questo tipo mostrano \gls{tagBias} \cite{hashimshony2012cel}.

Nel 2012, l'invenzione del \gls{fullLength} ha rimosso quel bias \cite{ramskold2012full, picelli2013smart}. Una molecola intera poteva ora essere copiata e letta da un capo all'altro, preservando la struttura interna che distingue una variante di trascritto dall'altra. Sebbene questo rappresentasse un guadagno sostanziale in copertura, la tecnologia era costosa e il formato su piastra manteneva uno studio nell'ordine di centinaia o poche migliaia di cellule.

Un nuovo protocollo, la microfluidica a goccia, ha rivoluzionato improvvisamente la scala nel 2015 \cite{zheng2017massively}: ogni goccia incapsula una cellula insieme a una sfera\defn{Sfera}{Una piccola sfera di gel caricata con i primer a codice a barre. Una sfera viene co-incapsulata con ogni cellula e fornisce a ciascuna goccia i primer che le servono.} che porta primer a codice a barre\defn{Primer a codice a barre}{Primer che portano una breve sequenza identificativa, così che ogni molecola catturata in una data goccia possa essere ricondotta alla sua cellula dopo il sequenziamento.}, quindi decine di migliaia di cellule vengono profilate in una sola corsa a una frazione del costo per cellula. Il guadagno ha avuto un prezzo. La cattura a goccia è ancorata alla coda del trascritto; questo impone l'uso della vecchia chimica, e il tag bias è tornato.

I \gls{dropletProtocols} dominano oggi il lavoro a singola cellula, perché profilano molte più cellule per esperimento a un costo per cellula molto inferiore rispetto a quanto consenta la chimica su piastra \cite{conte2024opportunities}. Ciò che il tag bias rimuove è quindi il limite centrale del formato, non un aspetto accessorio, e recuperarlo è una linea di ricerca attiva. Le piattaforme a lettura lunga leggono interi trascritti a risoluzione di singola cellula, e i dataset a lettura corta vengono sfruttati per qualunque segnale di isoforma i loro frammenti terminali ancora conservino \cite{joglekar2023words}.

\section{La sfida computazionale}
Due ostacoli si frappongono tra i dati a goccia e i risultati a livello di trascritto.

\textbf{\gls{sparsity} estrema.} Le \glspl{countMatrix} provenienti da esperimenti a goccia sono per lo più zeri, e le voci nulle costituiscono abitualmente dal 70\% al 90\% di tutti i valori. Alcuni di quegli zeri sono biologici e registrano un trascritto che la cellula non ha mai prodotto, mentre altri sono tecnici e registrano una molecola che era presente ma non è mai stata catturata, ed entrambi appaiono identici nella matrice.

\textbf{Read ambigui.} Di ogni molecola viene letto solo l'\gls{threePrimeEnd} terminale, quindi quando più \glspl{isoform} condividono quell'estremità appaiono identiche nei dati e le loro proporzioni non possono essere misurate direttamente.

Un sequenziamento più profondo riduce il primo ostacolo e non intacca il secondo: read aggiuntivi riempiono maggiormente la matrice, ma nessuna quantità di frammenti terminali recupera l'interno di un trascritto. Un metodo che funzioni sui dati a goccia deve quindi estrarre segnale da una matrice osservata solo in parte, accettando che la maggior parte di ogni molecola non sia mai stata vista.

\subsection*{Nota su ciò che resta recuperabile}

L'informazione a livello genico sopravvive al vincolo del \gls{tagBias}. Un gene viene registrato quando un read gli viene attribuito con sicurezza, e nella maggior parte dei casi un read che non può essere ricondotto a un singolo trascritto mappa comunque su isoforme dello stesso gene. Una matrice genica può quindi essere costruita dagli stessi read e resta informativa per l'analisi dell'\gls{deg}. Una proprietà conseguente è che questa matrice è più densa di quella dei trascritti.

\section{Approcci attuali e loro limiti}
Tre famiglie di metodi affrontano la regolazione a livello di trascritto, e ciascuna si ferma prima di rilevare la DTU su un intero dataset a goccia.

\textbf{Algoritmi bulk.} Strumenti consolidati per il \gls{bulkRnaSeq} come rMATS \cite{shen2014rmats}, SUPPA2 \cite{trincado2018suppa}, DRIMSeq \cite{nowicka2016drimseq} e DEXSeq \cite{anders2012dexseq} modellano proporzioni continue di isoforme a partire da una copertura profonda che attraversa le giunzioni. Applicati a matrici sparse a goccia, quegli assunti crollano e le stime risultanti sono instabili.

\textbf{Aggregazione in pseudo-bulk.} Una via d'uscita dal problema della copertura è l'\gls{pseudoBulk}, che somma le cellule di un cluster o di una condizione in un unico profilo per recuperare la profondità che gli strumenti bulk si aspettano \cite{gilis2021satuRn, tekath2021dturtle}. Il costo è la variazione cellula-cellula che il saggio è stato costruito per misurare: uno switch confinato a una piccola sottopopolazione viene diluito dalle cellule circostanti.

\textbf{Strumenti mirati per dati \gls{threePrime}.} Sierra \cite{patrick2020sierra} analizza direttamente i dati a goccia, ma solo per una classe ristretta di eventi alle estremità dei trascritti. Non offre alcun quadro genome-wide per lo \gls{isoformSwitch} multiplo a partire da conteggi grezzi ricchi di zeri.

Ciò che manca è un quadro che funzioni su matrici sparse a goccia a risoluzione di singola cellula, senza collassare le cellule tra loro e senza read a lunghezza intera.

\section{Obiettivo della tesi}

La domanda che questa tesi persegue nasce da COTAN (Co-expression Tables Analysis) \cite{galfre2021cotan}, un quadro per \glspl{countMatrix} a livello genico che concentra il proprio modello statistico sugli zeri della matrice anziché sui conteggi positivi. COTAN chiede se due geni vengono rilevati insieme meno spesso di quanto il caso preveda, usando tabelle di contingenza di presenza e assenza, quindi non richiede alcuna imputazione\defn{Imputazione}{Sostituire i conteggi nulli con valori previsti dalle cellule vicine, col rischio di inventare correlazioni spurie.} artificiale né alcuna trasformazione stabilizzante la varianza\defn{Trasformazione stabilizzante la varianza}{Una trasformazione applicata ai conteggi, come il logaritmo, affinché la loro varianza non dipenda più dalla loro media. Cambia la scala dei dati, trasformando una molecola non rilevata in un piccolo valore positivo anziché in uno zero.}, e si comporta bene sulle matrici sparse tipiche dei dati a singola cellula. Questa tesi chiede se le stesse statistiche reggano quando vengono applicate un livello più in basso, a matrici di \glspl{transcript} i cui conteggi sono ancora più radi: se conservino abbastanza informazione di \gls{isoform} da rilevare la \gls{dtu} senza risolvere il problema del \gls{tagBias}. L'obiettivo è volutamente modesto --- le isoforme che condividono la regione terminale restano indistinguibili --- e mira a mappare dove statistiche robuste alla sparsità recuperino comunque l'uso differenziale dei trascritti.

Il lavoro si è svolto dai read grezzi a un risultato riportato in quattro fasi, e il contributo della tesi è la strumentazione e l'analisi in ciascuna fase.

\textbf{Dai read alle matrici di trascritti.} Il punto di partenza è l'output grezzo di un esperimento a goccia, che la strumentazione standard normalmente collassa a conteggi a livello genico. Una pipeline Nextflow personalizzata è stata costruita in questa tesi per automatizzare l'intero percorso verso una matrice di conteggi: scarica i read, costruisce l'indice e richiama il workflow scrnaseq di nf-core per quantificare la matrice. La costruzione dell'indice offre due modalità. Una costruisce l'indice convenzionale a livello genico; l'altra costruisce un indice a livello di trascritto seguendo l'approccio che gli autori di simpleaf raccomandano per le matrici di trascritti, il quale mantiene separate le righe dei trascritti. Questo artefatto è la base di tutto ciò che segue, poiché l'analisi DTU richiede una risoluzione a livello di trascritto che le matrici geniche non possiedono.

\textbf{Una libreria di analisi.} L'analisi stessa è stata consolidata in una libreria R anziché in una serie di script ad hoc. La libreria fornisce uno strato di logging costruito sopra i log di Seurat e COTAN, così un'esecuzione riporta il proprio avanzamento e i propri errori in un'unica traccia coerente, ed espone funzioni che coprono l'intero workflow: leggere una matrice di conteggi e adattarla alla forma richiesta da COTAN, guidare i passi di analisi di COTAN, eseguire il rilevamento DTU personalizzato e confrontare le esecuzioni tra loro. Questa libreria è il secondo contributo, ed è dove le statistiche di COTAN, pensate per i geni, vengono trasformate in un rilevatore a livello di trascritto.

\textbf{I driver di analisi.} Uno strato di driver mette insieme le funzioni della libreria nel workflow completo, una configurazione per dataset e livello di feature, così che un'analisi completa venga riprodotta da un unico file dichiarativo anziché da comandi interattivi. Qui si uniscono la pipeline e la libreria, ed è questo che rende l'intero studio riproducibile da capo a fondo.

\textbf{Dai candidati ai risultati riportati.} Il workflow completo è stato applicato a due dataset --- una miscela omogenea di linee cellulari umane e un tessuto corticale murino complesso --- e gli switch candidati sono riportati in questa tesi, confrontati con la biologia delle isoforme pubblicata e con le baseline di clustering a livello genico. L'analisi è stata condotta in prima persona e non revisionata da un esperto del dominio, ed è presentata di conseguenza: i risultati sono offerti come un primo report riproducibile anziché come biologia consolidata, e ci si attende una verifica indipendente dei candidati prima di costruirvi sopra.
